\documentclass[10pt,a4paper]{amsart}
\usepackage[margin=19mm]{geometry}
\usepackage{amsmath,amssymb,mathtools}
\usepackage[scr=rsfs,cal=euler]{mathalfa}
\usepackage{xfrac}
\usepackage[unicode,hidelinks,bookmarks=false]{hyperref}
\newcommand{\ie}{i.e.\ }
\newcommand{\cf}{cf.\ }
\newcommand{\resp}{respectively\ }
\newcommand{\Subsection}{\subsection}
\providecommand{\nobreakdash}{\nobreak\hbox{-}}
\setlength{\emergencystretch}{2em}
\multlinegap0pt
\usepackage{amssymb}
\usepackage{bbm}
\usepackage{xcolor}
\usepackage[shortlabels]{enumitem}
\newtheorem{prop}{Proposition}
\newtheorem{thm}{Theorem}
\newtheorem{cor}{Corollary}
\newtheorem{lem}{Lemma}
\newcommand{\R}{\mathbb{R}}     % reelle Zahlen
\newcommand{\cN}{\mathcal{N}}  
\title{Optimal eigenvalues on a metric graph with densities}
\author{Kiyan Naderi, Noema Nicolussi}
\date{Source: arXiv:2512.20508v1 (CC BY 4.0)}
\begin{document}
\maketitle

\noindent\emph{Source and licence.} Optimal eigenvalues on a metric graph with densities. Version arXiv:2512.20508v1 (CC BY 4.0), distributed by arXiv under the Creative Commons Attribution 4.0 International licence, http://creativecommons.org/licenses/by/4.0/, as recorded in the arXiv licence metadata for this exact version and shown on the abstract page https://arxiv.org/abs/2512.20508v1. Full licence notice: \texttt{licenses/arxiv-2512.20508-CC-BY-4.0.txt}.\par\smallskip

\noindent\textit{Editorial scope.} Counted unit: the article's thm vertices (source0.tex lines 1397--1399). The article's complete original proof of it is delivered with it (source0.tex lines 1402--1453, one \texttt{proof} environment of 52 lines). No mathematics is written, repaired or re-derived: the statement and the proof are the article's own lines, in the article's own notation, and no retained line is substituted or re-flowed.  Retained uncounted as the source-local context the proof needs: lines 359--384; lines 464--471; lines 906--916; lines 918--924; lines 1050--1055; lines 1460--1462.  Nothing was omitted from any retained span, so the package carries no omission record.  No retained line contains a citation, so the wrapper carries no bibliography and no bibliography entry.  No retained line contains a figure environment, an image-inclusion command or drawing code, so no image is delivered and the package declares no asset dependency.  Editorial formatting only, in addition: an AMS \texttt{amsart} preamble that restates the article's own declarations and macros used by the retained lines, namely the environments \texttt{prop}, \texttt{thm}, \texttt{cor}, \texttt{lem} on separate counters, together with the standard packages those lines need.  The retained environments print in this document as Proposition 1, Theorem 1, Corollary 1, Lemma 1 in order of appearance, whereas the article prints the same items with the numbering of its own full document.  The complete native source of the article as distributed in the arXiv e-print for this exact version is bundled unchanged as \texttt{sources/191573/source0.tex}.
\begin{prop} \label{Prop: Harmonic extension}
	Let $\varnothing \neq A \subset V$ be finite and let $(f_x)_{x \in A} \subset \R$ be given function values.
	\begin{enumerate}
		\item [(i)] \label{1:HarmonicExtension} There exists a unique function $f \in H^1(G)$ such that 
		\begin{equation} \label{harmonic extension}
			\begin{cases}
			f(x) = f_x \text{ for all } x \in A, \\
			f \text{ is harmonic on } G\setminus A. 
		\end{cases}
		\end{equation}
	
	\item [(ii)] \label{2:HarmonicExtension} The solution of~\eqref{harmonic extension} is the unique function $f \in H^1(G)$ such that
		\begin{equation} \label{harmonic extension minimum}
			\begin{cases}
			f(x) = f_x, & \text{ for all } x \in A, \\
			Q(f)\le Q(g), & \text{ for all }  g \in \mathcal{F},
			\end{cases}
		\end{equation}
	where $\mathcal{F} = \{g \in H^1(G)| \, g(x) = f_x \text{ for all } x \in A\}$. 
	\item [(iii)] \label{3:HarmonicExtension} Let $f\in H^1(G)$ be harmonic on $G \setminus A$ and $g \in H^1(G)$. Then 
	\begin{equation} \label{harmonic extension IBP}
	q(g,f) =  \sum_{x \in A}g(x) \sum_{e: e \sim x } \frac{f(x) - f(y_e)}{\ell(e)},
	\end{equation}
	where $y_e \in V$ is the unique vertex which is connected to $x$ via $e$.  
	\end{enumerate}
\end{prop}

\begin{prop}[Parallel series for pumpkin subgraphs] \label{prop: surgery}
	Let $G=(V,E, \ell)$ be a metric graph and $f \in H^1(G)$ harmonic on an open subset $O\subset G$.
	Suppose $ \overline{O}$ contains two vertices $x_1, x_2$ and $n$ edges $e_1, \dots e_n \subset O$ between $x_1 $and $x_2$.

	Let $G'$ be the graph obtained by replacing the $n$ edges by a new edge $e$ connecting $x_1$ and $x_2$ with length $\ell'(e)=\left(\sum_{i=1}^{n} \ell(e_i)^{-1}\right)^{-1}$. Denote by $\Phi\colon G \to G'$ the natural surjective map, obtained by linearly dilating every edge $e_i$, $i = 1, \dots, n$, onto $e$.
	
	Then there exists a unique function $g$ on $G'$ with $g \circ \Phi = f$. Moreover, $g$ is harmonic on $\Phi(O)$ and $q(g) = q(f)$. 
\end{prop}

\begin{thm} \label{thm: characerisation of lambda1}
	Let $G$ be a metric graph. Then there exists a measure $\mu$ of the form
	\begin{equation} \label{eq:FormMinimizingMeasure}
		\mu = \frac{1}{2}\delta_x + \frac{1}{2} \delta_y
	\end{equation}
	for two points $x\neq y \in G$ such that
	\[
	\lambda_1^{\min}(G) = \lambda_1(H_\mu).
	\]
	Moreover, every measure $\mu\in \cN$ with $\lambda_1^{\min}(G) = \lambda_1(H_\mu)$ is of the form~\eqref{eq:FormMinimizingMeasure} for some points $x\neq  y \in G$.
\end{thm}

\begin{equation} \label{eq:Definitionfxy}
	\begin{cases}
		f(x)=1,    \\
		f(y) = -1, \\
		f \text{ harmonic}  \text{ on }G\setminus\{x,y\}.
	\end{cases}
\end{equation}

\begin{cor}\label{Cor: Alternative charakterization of lambda 1}
	Let $G$ be a metric graph. Then
	\[
	\lambda_{1}^{\min}(G) = \min_{\substack{f \in H^1(G) \\ \max f =1, \, \min f= -1}}q(f).
	\]
\end{cor}

\begin{lem} \label{lem:DecreasingEdgeLength}
	Let $G =(V,E, \ell)$ be a metric graph. Shortening an edge $e \in E$, that is, decreasing the length $\ell(e)$, increases the first optimal eigenvalue $\lambda_1^{\min}$. Moreover, contracting an edge $e \in E$, that is, collapsing it into one vertex, also increases $\lambda_1^{\min}$.
\end{lem}

\begin{thm}\label{thm: points avoid vertices}
	Let $\mu \in \cN$ be a measure with $\lambda_{1}^{\min}(G)=\lambda_1(H_{\mu})$ and write $\mu$ as $\mu = \mu_{x,y} = (\delta_x + \delta_y)/2$ for two points $x, y \in G$. Then $\deg(x), \deg(y) \in \{1,2\}$, i.e.~the two points lie at the end of a pending edge or in the interior of an edge.  
\end{thm}

\begin{proof}
	By Theorem~\ref{thm: characerisation of lambda1}, the minimizing measure $\mu$ is of the form $\mu = \mu_{x,y}$ for some points $x, y\in G$. Upon refining the model, we may assume that $x$ and $y$ are vertices.
	
	
	Assume that $\deg(x) > 2$ and let $f=f_{x,y}$. First contract every edge, where $f$ is constant. The new graph $G'$ has a bigger first optimal eigenvalue by Lemma~\ref{lem:DecreasingEdgeLength}. But we have $q(f) = q(f|_{G'})$, so the first optimal eigenvalue in fact coincides by Corollary~\ref{Cor: Alternative charakterization of lambda 1}. This also shows that $\mu_{x,y}$ is a minimizing measure for $G'$. Thus, we may assume that $G=G'$, that is, $f$ is non-constant on every edge.
	
	Since the function $f$ is linear and non-constant on edges, all its level sets $\{z \in G\; | \; f(z)=c\}$, $ c \in \R$, are finite. Denote the values, which are taken by $f$ in vertices, by $a_1 = 1 = f(x) > a_2 > \dots > a_n = -1 = f(y)$. We assume that for each $i$, the level set $V_i := \{z \in G | f(z) = a_i\}$ consists only of vertices. This can always be achieved by refining the model. In particular, $V = \bigsqcup_i V_i$.
	

	
	We construct a new graph $\widetilde G$ by glueing together all vertices $v \in V_i$ for each set $V_i$. Then $\widetilde G$ is a pumpkin chain with vertices $\widetilde w_1, \dots \widetilde w_n$ corresponding to the values $a_1, \dots a_n$. The points $x$ and $y$ can be identified with the endpoints $\widetilde w_1$ and $\widetilde w_n$ of $\widetilde G$, so that the pumpkin chain $\widetilde G$ starts with a $\deg(x)$-pumpkin and ends with a $\deg(y)$-pumpkin.
	
	
	
	The glueing induces a natural surjective map $\Phi\colon G \to \widetilde{G}$. By construction, the function $f$ is the pullback $f = \widetilde{ f} \circ \Phi $ of a function $\widetilde f \in H^1(\widetilde G)$ with $q(f) = q(\widetilde{ f})$. On the other hand, for every function $g \in H^1(\widetilde{G})$, the pullback $g \circ \Phi \in H^1(G)$ has the same energy, maximum and minimum. By Corollary~\ref{Cor: Alternative charakterization of lambda 1}, we infer
	\[
	q(\widetilde{ f})	\geq \lambda_{1}^{\min}(\widetilde{G}) \geq \lambda_{1}^{\min}(G) = q(f) = q(\widetilde{ f}).
	\]
	Applying Proposition \ref{prop: surgery}, we can collapse every pumpkin in $\widetilde G$ except the first one into a path of suitable length, to obtain a graph $\widehat{G}$ with the following property: for the corresponding surjective mapping $\Psi\colon \widetilde{G} \to \widehat{G}$, there exists a unique function $\widehat{f}$ with $\widehat{f} \circ \Psi = \widetilde{ f}$ and $q(\widehat{f}) = q(\widetilde{ f})$. The same argument as above implies that
	\[
	\lambda_{1}^{\min}(\widehat{G}) = q(\widehat{f}). 
	\]
	
	We proceed by calculating the energy $q(\widehat{f})$. The graph $\widehat{G}$ contains a $\deg(x)$-pumpkin, having $x$ as its left endpoint and a vertex $z$ as its right endpoint. Either $z=y$, and $\widehat G$ consist only of the pumpkin, or it contains an additional edge $e_y$ attached at $z$, having $y$ as its other endpoint. Let $\ell_1\geq \dots\geq \ell_{\deg (x)} >0$ be the lengths of the edges $e_1, \dots, e_{\deg(x)}$ in the pumpkin at $x$ and let $\ell_0$ be the length of the remaining graph. If the edge $e_y$ is absent, we have $\ell_0 = 0$. Further set
	\[
	S := \sum_{i=2}^{\deg(x)} \frac{1}{\ell_i}.
	\]
	In terms of these quantities, the energy of $\widehat f$ equals
	\[
	q(\widehat{f}) =  4 \, \frac{S + \frac{1}{\ell_1}}{\ell_0S + \frac{\ell_0}{\ell_1} +1}.
	\] 
	This follows from the fact that, by Corollary~\ref{Cor: Alternative charakterization of lambda 1} and Proposition~\ref{Prop: Harmonic extension}, the function $\widehat f$ is harmonic on $\widehat G\setminus \{x,y\}$ with $f(x) = 1$, $f(y) =-1$, and thus takes the values
	\[
	f(x) = 1, \qquad f(y) = - 1, \qquad  f(z) = \frac{\ell_0 S+ \frac{\ell_0}{\ell_1} - 1}{\ell_0 S + \frac{\ell_0}{\ell_1} +1},
	\]
	in the vertices $x,y$ and $z$ of $\widehat G$. 
	
	Using the assumption that $\deg(x) \ge 3$, we then construct a function $\widehat g \in H^1(\widehat G)$ with smaller energy than $\widehat f$. Note first that
	\begin{equation} \label{eq:DefinitionMysteryPoint}
		\ell_1S = \frac{\ell_1}{\ell_2} + \dots + \frac{\ell_1}{\ell_{\deg(x)}} \geq \deg(x) -1 \geq 2. 
	\end{equation}
	This permits us to define a point in the interior of the edge $e_1=[0, \ell_1]$ emanating at $x$ as $x' := \frac{\ell_1S  -1}{2S}$. 
	Let $\widehat g := f_{x', y} \in H^1(\widehat G)$ be the harmonic function on $\widehat G \setminus \{x',y\}$ defined in \eqref{eq:Definitionfxy}. We use Proposition \ref{prop: surgery} to calculate the energy of $\widehat{g}$. First collapse the edges $e_2, \dots, e_{\deg x}$ to a new edge of length $S^{-1}$. Then the new graph is a $[2,1]$-pumpkin chain between the vertices $x', z$ and $y$. By again collapsing the $2$-pumpkin, we get a new edge of length $\frac{\ell_1S+1}{4S}$. The resulting graph is a path graph of length $\frac{\ell_1S+1}{4S}+\ell_0$ and thus $\widehat g$ has energy
	\[
	q(\widehat g) = \frac{16S}{(4\ell_0 + \ell_1)S +1}.
	\]
	Using that $\ell_1 S \neq 1$ by~\eqref{eq:DefinitionMysteryPoint}, it is straightforward to deduce that
	\[
	q(\widehat g) <  q(\widehat f) =  \lambda_{1}^{\min}(\widehat{G}),
	\]
	which is a contradiction by Corollary~\ref{Cor: Alternative charakterization of lambda 1}.
\end{proof}

\end{document}
